\documentclass{article}
\usepackage[T2A]{fontenc}
\usepackage[utf8]{inputenc}   
\usepackage[english, russian]{babel}

% Set page size and margins
% Replace `letterpaper' with`a4paper' for UK/EU standard size
\usepackage[a4paper,top=2cm,bottom=2cm,left=2cm,right=2cm,marginparwidth=1.75cm]{geometry}

\usepackage{amsmath}
\usepackage{graphicx}
\usepackage[colorlinks=true, allcolors=blue]{hyperref}
\usepackage{amsfonts}
\usepackage{amssymb}
% \usepackage[left=1cm,right=1cm,top=1cm,bottom=1cm]{geometry}
\usepackage{hyperref}
\usepackage{seqsplit}
\usepackage[dvipsnames]{xcolor}
\usepackage{enumitem}
\usepackage{algorithm}
\usepackage{algpseudocode}
\usepackage{algorithmicx}
\usepackage{mathalfa}
\usepackage{mathrsfs}
\usepackage{dsfont}
\usepackage{caption,subcaption}
\usepackage{wrapfig}
\usepackage[stable]{footmisc}
\usepackage{indentfirst}
\usepackage{rotating}
\usepackage{pdflscape}
\usepackage{listings}

\usepackage{minted}

\usepackage{MnSymbol,wasysym}

\title{Домашнее задание 9}
\author{Лиза Фоменко}
\date{}

\begin{document}

\maketitle

\section*{\textbf{Задание 1}} \textbf{\textit{Вам нужно выбраться из лабиринта. Вы не знаете, сколько в нем комнат, и какая у него карта. По всем коридорам можно свободно перемещаться в обе стороны, все комнаты и коридоры выглядят одинаково (комнаты могут отличаться только количеством коридоров). Пусть $m$ - суммарное количество коридоров между комнатами. Предложите алгоритм, который находит выход из лабиринта или доказывает, что его нет, за $O(m)$ переходов между комнатами. В вашем расположении имеется неограниченное количество монет, которые вы можете оставлять в комнатах. Минотавр мертв, так что в лабиринте больше никого.}}

Мне кажется эту задачу можно легко решить, представив, что комнаты = вершины, а коридоры = ребра некоторого графа $G(V,E)$, где $|E| = m$. Тогда, предположим, что мы стоим в какой-то комнате из $v \in V$. Начнем обход графа-лабиринта в глубину, при этом мы бы хотели помечать пару комната-коридор. Так как у нас монеток не ограничено, то ничего не запрещает нам класть монетку в комнате около каждого входа в коридор, в который мы идем. Тогда алгоритм такой: 

\begin{verbatim}

first_v = v       # вершина, в которой мы находимся

while True:
    if curr_v == 'EXIT':            # если нашли выход
        return 'found exit!'
        
    elif в вершине-комнате curr_v есть непосещенные коридоры:    # если есть возможность продолжать DFS
        выбираем коридор
        кладем монетку около его входа
        идем по коридору в вершину-комнату new_v
        prev_v = curr_v
        curr_v = new_v
        
    else:                   # нашли комнату, где прошли все коридоры, возвращаемся обратно
        curr_v = prev_v

if curr_v == first_v:           # если вернулись туда, откуда начинали, значит, выхода нет
    return 'no exit found'
\end{verbatim}

Получается, что мы реализовали DFS, при этом каждую комнату мы посетили не более 2k раз, где k = количество коридоров, ведущих из этой комнаты, а каждое ребро - не более 2 раз (1 раз в одном направлении, 2 раз - в обратном). Как бы единственное отличие задачи от DFS в том, что мы используем монетки для мечения пройденных коридоров из данной комнаты. Коректность DFS мы вроде доказывали, и так как каждое ребро мы смотрим не более 2 раз, то сложность $O(m)$ (на самом то деле  $O(m + n)$ ну в алгоритмическом смысле $O(m+n)$, где $n$ = количество вершин-комнат. Но условно как бы ок, $O(m)$, потому что нам нужна асимптотика от переходов между комнатами, поэтому $O(m)$ works))))

\bigskip

\section*{\textbf{Задание 2}} \textbf{\textit{Предложите алгоритм, выясняющий, есть ли в неориентированном графе циклы нечетной длины. Оцените асимптотику его работы. Докажите его корректность.}}

Мы на лекции доказали, что любой ациклический граф двудольный. 

Докажем, что графы, имеющие циклы нечетной длины, не являются двудольными, а графы только с четными циклами двцдольными являются. Тогда нам останется только проверить графы на двудольность 

Докажем с помощью раскрашивания вершин.

Будем использовать обход в глубину. Давайте возьмем произвольную вершину $v_0$ и поставим ей метку 0. Для всех вершин на расстоянии 1 от нее ставим метку 1, для вершин на расстоянии 2 - метку 0 и так далее: для вершин на четном расстоянии ставим метку 0, для вершин на нечетном - 1.

Пусть теперь в нашем графе есть циклы, то есть существует путь из $v_i$ в $v_i$, проходящий через набор вершин $v_{i+1}, ..., v_{i+k}$. Пусть $v_i$ имеет метку 0. Для этой вершины так же характерно, что все вершины на четном расстоянии от нее окрашены так же, как она, то есть имеют метку 0, остальные - метку 1. Тогда, чтобы граф оставался двудольным, после прохождения по циклу, следуя этому правилу, натолкнувшись снова на вершину $v_i$, исходя из нее, мы должны вновь поставить ей метку 0, то есть "вершина $v_i$ находится на четном расстояния от себя самой", то есть длина ее цикла четна. Если же цикл будет нечетным, то присуждая четным и нечетным вершинам разыне метки, мы будем вынуждены поставить  $v_i$, имеющей метку 0, новую метку 1, то есть нашли противоречие, и наш граф двудольным не является. 

Доказали: если в графе есть цикл нечетной длины, значит, он не двудольный.

Дальше докажем обратное: если граф двудольный, то в нем нет циклов нечетной длины. Докажем. Пусть существует двудольный граф, который мы можем разделить на доли A, B, а в нем есть цикл нечетной длины из вершины $v_i$. Пусть $v_i$ принадлежит A. Тогда в B содержатся все вершины, находящиеся на нечетном расстоянии от $v_i$, а в A - на четном расстоянии. Мы знаем, что существует цикл нечетной длины $v_i$ -> $v_i$. Значит, существует путь нечетной длины от $v_i$ в $v_i$, значит, $v_i$ находится на нечетном расстоянии от самой себя и должна входить в долю B тоже. Но это невозможно по условию, так как в двудольном графе каждая вершина принадлежит только одной доле. Значит, предположение неверно.

Тогда мы знаем: любой ациклический граф двудольный (лекция); граф с циклами четной длины - двудольный (выше); граф с циклом нечетной длины - НЕ двудольный. Тогда проверка проста: проверям, является ли наш неорграф двудольным. Если нет, то это значит, что в нем есть цикл нечетной длины. Такую проверку можно делать путем окраски графа в 2 цвета с помощью DFS, который имеет асимптотику $O(|V| + |E|)$

\bigskip


\section*{\textbf{Задание 3}} \textbf{\textit{Предложите $O(|V| + |E|)$ алгоритм поиска максимального по размеру независимого множества в дереве.}}

Так как нам дано дерево - связный ациклический граф, то здесь независимое множество выглядит примерно так: выберем к-л вершину $v_0$ в качестве корня этого дерева. 
Тогда смежными для нее будут являться только ее дети, а их дети смежными являться уже не будут. Аналогично для каждой вершины жтого дерева: она является смежной только для своих прямых детей и для своего прямого родителя (существует единственный родитель этой вершины, иначе мы бы получили цикл!). Тогда просто разделив все вершины по четности/нечетности удаления от корня дерево мы получим два независимых множества в этом дереве.

Но мы ищем максимальное, поэтому просто заметим, что для любой вершины верно, что в максимальное независимое множество может входить либо сама эта вершина, либо ее потомки: вместе не могут. То есть на этапе рассмотрения вершины нам необходимо уметь обрабатывать 2 случая: вершина $v$ входит в множетсво => ее дети не могут входить в множетсво; вершина $v$ не входит в множество => ее дети могут в него входить. Также важно, что любая вершина дерева образует собственное поддерево, для которого верны все те же рассуждения, что и для целого дерева.

Обычно подобные задачи решаются динпрогом: делим наше дерево на поддеревья, для каждого из них решаем свою задачу, и комбинируем решения. Также по асимптотике я предполагаю, что нужен какой-то обход дерева

Напишем рекурсивный алгоритм:


\begin{verbatim}

def find_max_ind_set(curr_v):

    sum_when_isin = 1         # если curr_v входит в множество, то все ее дети точно не входят
    all_v_when_isin = [curr_v]      # список вершин, составляющих мн-во
    sum_when_absent = 0       # если curr_v не входят в множество, то дети могут как входить, так и не входить в него
    all_v_when_absent = []          # список вершин, составляющих мн-во
    

    for v' in curr_v_children:         # идем рекурсивно по всем детям дерева
        # находим максимальные размеры независимых множеств для поддеревьев с корнем в v'
        child_isin, child_absent, vs_child_isin, vs_child_absent = find_max_ind_set(v')  

        # теперь к sum_when_isin добавляем child_absent - так как вершина  
        # curr_v не может входить в множество вместе с детьми
        sum_when_isin += child_absent             
        all_v_when_isin += vs_child_absent

        
        # к sum_when_absent добавляем максимум из child_isin, child_absent, 
        # так как все потомки независимы друг от друга (граф ацикличен)
        if child_isin > child_absent:
            sum_when_absent += child_isin
            all_v_when_absent += vs_child_isin
        else:
            sum_when_absent += child_absent
            all_v_when_absent += vs_child_absent

    # функция возвращает максимальное по размеру независимое множество # (максимум из включающих и нет вершину curr_v)
    
    return sum_when_isin, sum_when_absent, all_v_when_isin, all_v_when_absent     
        

root = tree.root               # выбираем корень дерева
root_isin, root_absent, vs_root_isin, vs_root_absent = find_max_ind_set(root)    # запускаем из корня

if root_isin > root_absent:
    print(root_isin, vs_root_isin)      # выводим размер и входящие вершины
else:
    print(root_absent, vs_root_absent)      # выводим размер и входящие вершины
\end{verbatim}

За сколько работает алгоритм? за O(|E| + |V|), так как мы каждую вершину рассматриваем только один раз (то есть алгоритм запускается для нее единственный раз). Само тело алгоритма состоит из нескольких мат операций (сравнение и сложение), работающих за O(1), и мерджа списков вершин, которые не пересекаются (идем снизу вверх, у любой вершины существует единсвтенный родитель - иначе цикличность).

\bigskip


\section*{\textbf{Задание 4}} \textbf{\textit{На вход задачи поступает описание двудольного графа $G(L, R, E)$, степень каждой вершины которого равна двум. Необходимо найти максимальное паросочетание в $G$ (которое содержит максимальное количество рёбер). Предложите алгоритм, решающий задачу за $O(|V | + |E|)$.}}

Граф двудольный => либо ациклический, либо имеет только циклы четной длины (доказала выше). 

Рассмотрим просто цикл четной длины. Тривиально, что здесь максимальное паросочетание (их на самом деле два) имеет длину n/2, где n - длина цикла, и для его составления нужно просто взять каждое второе ребро.

Здесь не очень понятно как использовать информацию о степенях вершин. Если граф неориентированный, то рассмотрим к-л цикл из n ребер. Каждое ребро соединяет ровно 2 вершины графа, то есть добавляет +1 в степень каждой из них. При этом цикл подразумевает, что каждая вершина соединена с двумя другими. Тогда вывод: если в таком графе есть цикл, то он обязательно нечетный, а из условия на степени вершин вытекает, что все циклы изолированы - цикличный граф несвязный, каждый цикл изолирован. Значит, мы с помощью DFS находим все циклы, из них выбираем каждое второе ребро. Если общая длина циклов = n (четная, так как сумма четных чисел), то мы получим n/2 ребер.

Могут ли быть ациклические подграфы? Ответ - нет, так как любая вершина связана с двумя другими вершинами, а значит. Вообще граф, где степени всех вершин одинаковы, называются регулярными графами, и известно, что 2-регулярный граф состоит из разрозненных циклов. А для каждого цикла алгоритм мы уже нашли.

Тогда алгоритм прост: так как граф состоит из разрозненных циклов четной длины, то нам необходимо лишь сделать обычный обход в глубину DFS, а начинать, например, каждый раз с вершин из L. Считаем длину цикла n, а к размеру максимального паросочетания сохраняем n/2. Алгоритм займет O(|E| + |V|), как и обход DFS.
\bigskip


\section*{\textbf{Задание 5}} \textbf{\textit{На вход задачи поступает граф $G(V, E)$ в виде списка исходящих ребер для каждой вершины. Опишите алгоритм, переводящий это представление графа в представление в виде матрицы смежности...}}

\begin{verbatim}

GIVEN: # dictionary: v - merged
{a: [b, c, e]
b: [a, k, l]
....
}


WANTED:
|___|_a_|_b_| ...
|_a_| - | 1 | ...
|_b_| 1 | - | ...
...

ALGORITHM

all_v = {}
i = 0


for v, merged in merged_dict.items():       # for every mergence list
    if v not in all_v.keys():      # add all v in dict, assign unique index value
        all_v[v] = i               # O(1)
        i += 1                     # O(1)
    for v' in merged:
        if v' not in all_v.keys():
            all_v[v] = i           # O(1)
            i += 1                 # O(1)

# current total: O(|E|)
# now we know total number of V and index for every v
# we can make matrix

n = len(all_v.keys())               # |V|
matrix = [[0]*n]*n                  # |V|*|V| matrix 

for v, merged in merged_dict.items():       # for every mergence list
    v_index = all_v[v]                      # v index in matrix indexes

    for v' in merged:
        v'_index = all_v[v']                   # v' index in matrix indexes
        matrix[v_index, v'_index] = 1          # we have an edge here
        
TOTAL: O(|E|)
\end{verbatim}

Самым трудозатратным этапом будет сборка всех вершин: нам нужно найти все вершины, которые есть в данном графе. Замечу, что на граф не накладывается никаких условий, поэтому он может быть как орграфом, так и неорграфом: мы не можем взять ключи из словаря со списками смежности, так как может существовать вершина орграфа, в которую ребра только заходят, но не выходят из нее. Поэтому нам придется пройтись по каждому списку словаря смежности: мы это сделаем за $O(|E|)$: так как количество ключей оценивается через $|V|$, а $len(arr) = |e(v,u)|$, где $v$ - вершина из ключа, то общая длина списков равна $|E|$. 

Эти вершины мы будем класть в новый словарь вида $v: i$, где $v$ - вершина, а $i$ - ее порядковый номер. Добавление и извлекание из хеш-таблиц выполняется за $O(1)$. Эта операция нужна для того, чтобы конвертировать названия вершин в индексы нашей таблицы, от 0 до |V|-1. 

Теперь all\_v содержит все вершины, из которых исходят ребра, а также их индексы в нашей матрице смежности.

Создадим матрицу смежности. Пусть это выполняется за O(1), например, списком списков в питоне.

Теперь мы еще раз пройдемся по всем спискам смежности. Теперь мы можем по индексам вершин из all\_v добавить каждое ребро за $O(1)$ путем доступа по индексу, тогда всего O(|E|) операций.

Таким образом: 

- создать индекс для каждой вершины, из которой выходит ребро, можно за O(|E|)

- создать пустую матрицу смежности - условно O(1)

- заполнить матрицу смежности - за O(|E|)

total: O(|E|)
\bigskip

\section*{\textbf{Задание 6}} \textbf{\textit{...и наоборот. Докажите корректности и оцените сложность обоих алгоритмов.}}

Когда у нас дана матрица смежности, то она имеет размер |V|*|V|. При этом у нас названия вершин - если нужны конкретные, например, буквенные, могут быть заданы индексами самой матрицы, или лежат в отдельном словаре. Тут допустим, что операция доступа к имени вершины у нас осуществляется за O(1) путем индексирования или доступа к словарю индекс-название. 

\begin{verbatim}

GIVEN: 
|___|_a_|_b_| ...
|_a_| - | 1 | ...
|_b_| 1 | - | ...



WANTED:
...
{a: [b, c, e]
b: [a, k, l]
....
}

ALGORITHM

mergence_dict = {}
n = len(matrix)

for i in range(n):                     # |V| times
    v = names[i]                       # get v name, O(1)
    v_list = []
    for j in range(n):                 # |V| times
        if matrix[i,j] == 1:           # O(1)
            v_list.append(names[j])    # append v', O(1)
    mergence_dict[v] = v_list
        
TOTAL: O(|V|^2)
\end{verbatim}

Итого мы просматриваем каждую ячейку матрицы смежности и обрабатываем ее за O(1). Так как размер матрицы |V|*|V|, то сложность O(|V|$^2$)

\bigskip

\end{document}